\documentclass[11pt,a4paper]{article}

%─────────────────────────── Encoding / Fonts ──────────────────────────────
\usepackage[utf8]{inputenc}
\usepackage[T1]{fontenc}
\usepackage{lmodern}
\usepackage{microtype}

%─────────────────────────── Layout ────────────────────────────────────────
\usepackage[a4paper, top=1in, bottom=1.15in, left=1in, right=1in]{geometry}
\usepackage{setspace}
\usepackage{float}

%─────────────────────────── Math ──────────────────────────────────────────
\usepackage{amsmath,amssymb}

%─────────────────────────── Tables ────────────────────────────────────────
\usepackage{booktabs,longtable,array,tabularx,colortbl}

%─────────────────────────── Colors ────────────────────────────────────────
\usepackage{xcolor}
\definecolor{NavyBlue}{RGB}{0,58,120}
\definecolor{DarkGreen}{RGB}{0,110,60}
\definecolor{CritRed}{RGB}{175,30,30}
\definecolor{WarnAmber}{RGB}{160,90,0}
\definecolor{MidGray}{RGB}{200,200,210}
\definecolor{TableHead}{RGB}{225,233,245}
\definecolor{TableBorder}{RGB}{160,178,210}
\definecolor{CodeBg}{RGB}{247,248,252}
\definecolor{WarnYellow}{RGB}{255,249,220}
\definecolor{AlertRed}{RGB}{254,236,236}

%─────────────────────────── Code Listings ─────────────────────────────────
\usepackage{listings}
\lstdefinestyle{mypython}{
  language=Python,
  basicstyle=\ttfamily\small,
  keywordstyle=\bfseries\color{NavyBlue},
  commentstyle=\itshape\color{DarkGreen},
  stringstyle=\color{CritRed!80!black},
  showstringspaces=false,
  breaklines=true,
  breakatwhitespace=true,
  frame=single,
  framesep=4pt,
  rulecolor=\color{MidGray},
  backgroundcolor=\color{CodeBg},
  columns=fullflexible,
  xleftmargin=1.5em,
  xrightmargin=0.5em,
  aboveskip=0.8em,
  belowskip=0.8em,
  numbers=left,
  numberstyle=\tiny\color{gray},
  numbersep=7pt,
}

%─────────────────────────── Links ─────────────────────────────────────────
\usepackage{hyperref}
\hypersetup{
  colorlinks=true,
  linkcolor=NavyBlue,
  urlcolor=NavyBlue,
  citecolor=DarkGreen,
  pdftitle={MMHCL+ Revision43 Implementation Review},
  pdfauthor={Technical Review},
  pdfsubject={Neighbor-Layer Hypergraph Contrastive Learning}
}

%─────────────────────────── Symbols ───────────────────────────────────────
\usepackage{enumitem}
\usepackage{pifont}
\newcommand{\cmark}{\textcolor{DarkGreen}{\ding{52}}}
\newcommand{\xmark}{\textcolor{CritRed}{\ding{56}}}

%─────────────────────────── Section Headings ──────────────────────────────
\usepackage{titlesec}
\titleformat{\section}
  {\large\bfseries\color{NavyBlue}}{\thesection.}{0.6em}{}
  [\vspace{-2pt}{\color{NavyBlue}\rule{\linewidth}{0.5pt}}]
\titleformat{\subsection}
  {\normalsize\bfseries\color{NavyBlue!80!black}}{\thesubsection.}{0.5em}{}
\titlespacing{\section}{0pt}{1.6em}{0.6em}
\titlespacing{\subsection}{0pt}{1.2em}{0.3em}

%─────────────────────────── Callout Boxes ─────────────────────────────────
\usepackage[most]{tcolorbox}
\tcbuselibrary{skins,breakable}

\newtcolorbox{summarybox}{
  enhanced, breakable,
  colback=NavyBlue!5!white,
  colframe=NavyBlue,
  arc=4pt, boxrule=0.8pt,
  left=8pt, right=8pt, top=8pt, bottom=8pt,
  title={\bfseries Executive Summary},
  fonttitle=\bfseries\color{white},
  attach boxed title to top left={yshift=-2.5mm, xshift=5mm},
  boxed title style={colback=NavyBlue, arc=3pt, boxrule=0pt, left=4pt, right=4pt}
}

\newtcolorbox{alertbox}[1]{
  enhanced, breakable,
  colback=AlertRed,
  colframe=CritRed,
  arc=3pt, boxrule=0.7pt,
  left=6pt, right=6pt, top=5pt, bottom=5pt,
  fontupper=\small,
  title={\bfseries\small #1},
  coltitle=white,
  attach boxed title to top left={yshift=-2mm, xshift=5mm},
  boxed title style={colback=CritRed, arc=2pt, boxrule=0pt, left=3pt, right=3pt},
  before=\vspace{0.5em},
  after=\vspace{0.3em}
}

\newtcolorbox{warnbox}[1]{
  enhanced, breakable,
  colback=WarnYellow,
  colframe=WarnAmber,
  arc=3pt, boxrule=0.7pt,
  left=6pt, right=6pt, top=5pt, bottom=5pt,
  fontupper=\small,
  title={\bfseries\small #1},
  coltitle=WarnAmber!80!black,
  attach boxed title to top left={yshift=-2mm, xshift=5mm},
  boxed title style={colback=WarnYellow, colframe=WarnAmber, arc=2pt,
                     boxrule=0.6pt, left=3pt, right=3pt},
  before=\vspace{0.5em},
  after=\vspace{0.3em}
}

%─────────────────────────── Header / Footer ───────────────────────────────
\usepackage{fancyhdr}
\pagestyle{fancy}
\fancyhf{}
\fancyhead[L]{\small\textcolor{NavyBlue}{\textbf{MMHCL+ Revision\,43 --- Implementation Review}}}
\fancyhead[R]{\small\textcolor{gray}{April 10, 2026}}
\fancyfoot[C]{\small\textcolor{gray}{\thepage}}
\renewcommand{\headrulewidth}{0.4pt}
\renewcommand{\footrulewidth}{0pt}
\setlength{\headheight}{14pt}

%─────────────────────────── Spacing ───────────────────────────────────────
\setstretch{1.15}
\setlength{\parskip}{0.55em}
\setlength{\parindent}{0pt}
\setlength{\arrayrulewidth}{0.5pt}

%─────────────────────────── Title ─────────────────────────────────────────
\title{%
  \vspace{-0.8em}%
  {\Large\bfseries\color{NavyBlue} Implementation Review of Revision\,43}\\[0.5em]
  {\normalsize\color{gray}
    Neighbor-Layer Hypergraph Contrastive Learning $\longrightarrow$ MMHCL+}%
  \vspace{-0.5em}%
}
\author{}
\date{}

\begin{document}
\maketitle
\thispagestyle{fancy}
\vspace{-1.2em}

\begin{summarybox}
I carefully read the entire \TeX{} specification file (revision\,43) and cross-checked every component against the Python source code.
The system implements most of the design correctly; however, \textbf{five important deviations} from the \TeX{} description require attention.
Overall implementation completeness is estimated at \textbf{$\approx$\,75\%}.
The training loop, loss functions, and scaffold structure are already solid.
The three core novelty elements --- EMA teacher, real GradNorm, and FAISS-based Soft Topology-Aware Purification --- still need to be added to reach 100\%.
\end{summarybox}

\section{Components Implemented Correctly at 100\%}

{\renewcommand{\arraystretch}{1.35}%
\arrayrulecolor{TableBorder}%
\begin{longtable}{|>{\raggedright\arraybackslash}p{0.052\textwidth}
  |>{\raggedright\arraybackslash}p{0.435\textwidth}
  |>{\raggedright\arraybackslash}p{0.295\textwidth}
  |>{\raggedright\arraybackslash}p{0.12\textwidth}|}
\hline
\rowcolor{TableHead}
\textbf{\#} & \textbf{Component (\TeX{} Section)} & \textbf{Implementation File} & \textbf{Status} \\
\hline
\endfirsthead
\hline
\rowcolor{TableHead}
\textbf{\#} & \textbf{Component (\TeX{} Section)} & \textbf{Implementation File} & \textbf{Status} \\
\hline
\endhead
1 & Expanded Projector $d\!=\!64 \rightarrow 2048 \rightarrow 8192$, 3-layer MLP + BN + ReLU (\S4.4) & \texttt{projector.py}, \texttt{config.py} & \cmark\ Exact \\\hline
2 & Barlow Twins loss + optional \texttt{soft\_weights} (\S4.4) & \texttt{losses.py} --- \texttt{barlow\_twins\_loss()} & \cmark\ Exact \\\hline
3 & Chunked InfoNCE + \texttt{dynamic\_weights} + \texttt{chunk\_size} (\S4.4) & \texttt{losses.py} --- \texttt{chunked\_info\_nce\_loss()} & \cmark\ Exact \\\hline
4 & Soft BYOL alignment + stop-gradient (\S4.4 Algorithm~1, Step~6) & \texttt{soft\_byol.py} --- \texttt{.detach()} on target & \cmark\ Exact \\\hline
5 & Neighbor-layer pairs $(E^{(l)}, E^{(l+1)})$ + stop-grad (\S4.3) & \texttt{neighbor\_pairs.py} & \cmark\ Exact \\\hline
6 & Dirichlet Energy $\mathrm{tr}\!\left(E^\top(I-\Theta)E\right)/n$ (\S3.3) & \texttt{dirichlet.py} + inline \texttt{sparse\_dirichlet\_energy()} & \cmark\ Both correct \\\hline
7 & Temperature annealing $\tau = \max(\tau_{\min}, \tau_{\max} \cdot \gamma^t)$ (\S4.4) & epoch loop in \texttt{main\_mmhcl\_plus.py} & \cmark\ Exact \\\hline
8 & Warmup epochs --- BPR-only phase (\S4.4) & \texttt{main\_mmhcl\_plus.py} --- \texttt{in\_warmup} gate & \cmark\ Exact \\\hline
9 & 5-task loss aggregation (BPR, u2u, i2i, align, Dir) (Algorithm~1, Step~9) & \texttt{gradnorm.py} + \texttt{balancer.combine()} & \cmark\ All 5 present \\\hline
10 & \texttt{WEMAManager} --- lazy row fetch + periodic update (\S4.2) & \texttt{dynamic\_ema\_weights.py} & \cmark\ Correct \\\hline
11 & \texttt{MMHCLWithLayers} --- layer caching for ii and uu branches & \texttt{main\_mmhcl\_plus.py} --- \texttt{forward\_plus()} & \cmark\ Correct caching \\\hline
12 & Package-level extension, original MMHCL backbone preserved intact (\S4.1) & \texttt{mmhcl\_plus\_scaffold/} & \cmark\ Correct \\\hline
\end{longtable}%
\arrayrulecolor{black}}

\noindent The components above map one-to-one to Algorithm~1, Listings~1--2, and Figures~3--4 in the \TeX{} specification.

\section{Five Critical Deviations That Must Be Fixed}

\subsection{Deviation 1 --- Missing EMA Teacher Network (Algorithm~1, Steps 6 and~11)}

\begin{alertbox}{Critical \textemdash\ EMA teacher network is completely absent}
The target view is currently taken directly from the online bipartite embeddings,
which \textbf{violates the core BYOL design principle} described in the \TeX{} specification.
\end{alertbox}

The \TeX{} document states this clearly: Soft BYOL alignment must use an EMA teacher network $f_{\xi}$ --- a momentum-updated copy of the encoder --- to generate the target view, followed by applying $sg(\cdot)$ (stop-gradient).
This is the core idea inherited from BYOL (Grill et al., 2020).
Figure~4 in the \TeX{} specification explicitly shows an \emph{EMA Teacher Network (momentum\,=\,0.99)} as a separate processing stream.

\textbf{Current code behavior:} \texttt{main\_mmhcl\_plus.py} directly uses \texttt{i\_bip} / \texttt{u\_bip} (bipartite embeddings) with \texttt{.detach()}:

\begin{lstlisting}[style=mypython, caption={Current (incorrect) alignment call}]
batch_aln = (
    soft_byol_alignment(ii_final, i_bip)     # missing EMA teacher!
    + soft_byol_alignment(uu_final, u_bip)
)
\end{lstlisting}

\textbf{Required fix:} Add an EMA teacher model as follows.
\begin{itemize}[leftmargin=1.5em, itemsep=2pt]
\item \texttt{ema\_model = copy.deepcopy(model)}
\item a decorated function \texttt{@torch.no\_grad() update\_ema(student, teacher, momentum=0.99)}
\item invoke the EMA update at the end of every mini-batch (Algorithm~1, Step~11):
\end{itemize}

\[
  \xi \;\leftarrow\; \beta\,\xi \;+\; (1-\beta)\,\theta
\]

The target view must then be produced from \texttt{ema\_model.forward\_plus(\ldots)} rather than from the online bipartite embeddings directly.

\subsection{Deviation 2 --- ``GradNorm'' Is Only a Learned Softmax, Not Real GradNorm}

\begin{warnbox}{Warning \textemdash\ Completeness $\approx$ 30\%; only the weighting structure exists}
\texttt{self.alpha} is stored but never participates in any gradient-norm computation.
The module behaves as a simple learnable weighting layer.
\end{warnbox}

\TeX{} \S4.4 explicitly refers to a \emph{GradNorm Dynamic Balancing System}.
The actual GradNorm algorithm (Chen et al., 2018) requires:
\begin{enumerate}[leftmargin=1.5em, itemsep=2pt]
\item computing the gradient norm of each task loss with respect to shared parameters,
\item computing the relative training rate of each task,
\item performing a secondary update step on task weights to equalize those rates.
\end{enumerate}

\textbf{Current code in \texttt{gradnorm.py}:}

\begin{lstlisting}[style=mypython, caption={Current \texttt{combine()} \textendash{} only softmax weighting, no gradient norms}]
def combine(self, losses):
    weights = self.normalized_weights()   # only softmax(logits)
    return (weights * loss_vec).sum()     # alpha is never used
\end{lstlisting}

\texttt{self.alpha} is stored but never participates in the computation.
Therefore this module is effectively a \emph{learnable weighting layer}, not true GradNorm.
The \texttt{combine()} function must be rewritten to implement genuine gradient-norm equalization.

\subsection{Deviation 3 --- Missing FAISS/LSH and Soft Topology-Aware Purification (\S3.4)}

\begin{warnbox}{Warning \textemdash\ Completeness $\approx$ 20\%; this is the largest algorithmic gap}
Only dense cosine similarity is computed. The Jaccard term, bridge-score penalty,
and FAISS/LSH backend are entirely absent.
\end{warnbox}

\TeX{} \S3.4 describes the Soft Topology-Aware Purification mechanism in detail, using four signals:
\begin{itemize}[leftmargin=1.5em, itemsep=2pt]
\item cosine similarity,
\item sampled Jaccard coefficient,
\item popularity penalty (bridge score),
\item percentile score,
\end{itemize}
combined through the \texttt{soft\_topology\_weight()} formula and accelerated with ANN/FAISS so that complexity becomes $\mathcal{O}(N \cdot k)$ instead of $\mathcal{O}(N^2)$.

\textbf{Current code behavior:} \texttt{WEMAManager} uses only a dense cosine similarity matrix.
There is no Jaccard term, no bridge-score penalty, and no FAISS/LSH backend.
This is the largest deviation in terms of algorithmic novelty.

\subsection{Deviation 4 --- ASW Applied Only on the Item Side; User Side Missing}

\begin{alertbox}{Critical \textemdash\ \texttt{w\_ema\_u} for users does not exist}
\texttt{soft\_weights} is not passed when Barlow Twins is called for the u2u branch,
meaning the user-side branch is effectively unweighted.
\end{alertbox}

\TeX{} Listing~2 (\S4.4) clearly applies Adaptive Sample Weighting (ASW) to \textbf{both} branches:

\begin{lstlisting}[style=mypython, caption={Expected ASW usage from \TeX{} Listing~2}]
soft_weights    = FAISS_EMA_W['u2u'][batch_users]   # u2u branch
dynamic_weights = FAISS_EMA_W['i2i'][batch_items]   # i2i branch
\end{lstlisting}

\textbf{Current code behavior:}
\begin{itemize}[leftmargin=1.5em, itemsep=2pt]
\item only \texttt{self.w\_ema\_i} exists for items;
\item \texttt{w\_ema\_u} for users is absent;
\item \texttt{soft\_weights} is not passed when Barlow Twins is called for the u2u branch.
\end{itemize}

Current (incorrect) call pattern:

\begin{lstlisting}[style=mypython, caption={Current (incorrect) u2u branch call}]
u2u_terms.append(
    barlow_twins_loss(z1, z2, lambd=args.barlow_lambda)
    # missing: soft_weights = w_ema_u.get_batch_weights(users)
)
\end{lstlisting}

\subsection{Deviation 5 --- Profiling-Guided Checkpointing Not Wired into \texttt{forward\_plus()}}

\begin{warnbox}{Warning \textemdash\ Completeness $\approx$ 50\%; module exists but is never invoked}
\texttt{checkpointing.py} correctly implements the intended logic, but it is
not called anywhere inside \texttt{main\_mmhcl\_plus.py}.
\end{warnbox}

\TeX{} \S4.3 Table~2 specifies: \emph{``Checkpoint deep layers only; cache shallow layers ($g \leq 2$).''}
The file \texttt{checkpointing.py} already implements this logic correctly.

\textbf{Current code behavior:} \texttt{MMHCLWithLayers.forward\_plus()} calls \texttt{torch.sparse.mm()} directly without wrapping through checkpointing:

\begin{lstlisting}[style=mypython, caption={Current (uncheckpointed) forward pass}]
for _ in range(self._n_item_layers):
    ii_emb = torch.sparse.mm(I2I_mat, ii_emb)  # not checkpointed!
\end{lstlisting}

The module \texttt{profiling\_guided\_checkpoint} exists in the scaffold but is not invoked anywhere inside \texttt{main\_mmhcl\_plus.py}.

\section{Completion Summary}

{
\renewcommand{\arraystretch}{1.35}
\arrayrulecolor{TableBorder}%
\begin{longtable}{|>{\raggedright\arraybackslash}p{0.27\textwidth}
  |>{\raggedright\arraybackslash}p{0.14\textwidth}
  |>{\raggedright\arraybackslash}p{0.17\textwidth}
  |>{\raggedright\arraybackslash}p{0.32\textwidth}|}
\hline
\rowcolor{TableHead}
\textbf{Component} & \textbf{\TeX{} Section} & \textbf{Completion} & \textbf{Notes} \\
\hline
\endfirsthead
\hline
\rowcolor{TableHead}
\textbf{Component} & \textbf{\TeX{} Section} & \textbf{Completion} & \textbf{Notes} \\
\hline
\endhead
Package structure & \S4.1 & \cmark\ 100\% & Even better than the \TeX{} scaffold \\\hline
Expanded Projector & \S4.4 & \cmark\ 100\% & $d \rightarrow 2048 \rightarrow 8192$ is correct \\\hline
Barlow Twins & \S4.4 & \cmark\ 100\% & Both code and math are correct \\\hline
Chunked InfoNCE & \S4.4 & \cmark\ 100\% & Chunking + dynamic weights are correct \\\hline
Neighbor-layer pairs & \S4.3 & \cmark\ 95\% & Missing explicit \texttt{max\_hops} $= g$ limit \\\hline
Temperature annealing & \S4.4 & \cmark\ 100\% & Formula is correct \\\hline
Warmup epochs & \S4.4 & \cmark\ 100\% & BPR-only phase is correct \\\hline
Dirichlet energy & \S3.3 & \cmark\ 100\% & Both implementations are correct \\\hline
EMA Teacher Network & \S4.4 Alg.\,1 & \xmark\ 0\% & Completely missing \\\hline
GradNorm & \S4.4 & \textcolor{WarnAmber}{\textbf{30\%}} & Only learned weighting, not real GradNorm \\\hline
FAISS + Soft Purification & \S3.4, \S4.2 & \textcolor{WarnAmber}{\textbf{20\%}} & Only cosine similarity exists \\\hline
ASW for u2u branch & \S4.4 & \xmark\ 0\% & Implemented only on the item side \\\hline
Checkpoint in forward pass & \S4.3 & \textcolor{WarnAmber}{\textbf{50\%}} & Module exists but is not wired in \\\hline
\end{longtable}
\arrayrulecolor{black}%
}

\vspace{0.4em}
\noindent Overall, the system currently matches approximately \textbf{75\%} of the \TeX{} design.
The training loop, loss functions, and scaffold structure are already solid.
However, the three core elements that define the paper's novelty --- the EMA teacher, real GradNorm, and FAISS-based Soft Topology-Aware Purification --- must still be added in order to reach revision\,43 at 100\%.

\end{document}
